\section{The sharp all-core inequality at rank five}\label{sec:forcedcore}
The full lift gives a short general boundary theorem before the independent rank-five certificate below.
\begin{theorem}[Sharp all-core constant for a two-vertex shore]\label{thm:generalcore}
For a $2+q$ cover with $q\ge2$, factor the product of all core activities from the all-core sector and write it as $t^q(c_0+c_1t+c_2t^2)$. Then
\[
 c_1^2\ge\frac{2q}{q-1}c_0c_2,
\]
and the constant is optimal for every $q$.
\end{theorem}
\begin{proof}
In the lifted matroid basis polynomial, set all private $B$ variables to zero and differentiate once in each augmented core variable $X,Y$. These operations preserve the Lorentzian property, giving a degree-$q$ polynomial. Under the remaining specializations its three plane-count sectors give, up to the fixed positive scaling factor,
\[
 c_0s^2t^{q-2}+c_1st^{q-1}+c_2t^q.
\]
The bivariate $\ULC_q$ inequality at index $q-1$ is the desired bound. Zero activities and $Q=0$ follow by continuity as before.

For the complete $2\times q$ core with $N\ge q$ unit common-left vertices and $M\ge2$ unit common-right vertices,
\[
 c_0=\binom N{q-2},\qquad c_1=M\binom N{q-1},\qquad
 c_2=\binom M2\binom Nq.
\]
Their ratio is
\[
 \frac{c_1^2}{c_0c_2}
 =\frac{2q}{q-1}\frac{M}{M-1}\frac{N-q+2}{N-q+1}
 \longrightarrow\frac{2q}{q-1}.
\]
\end{proof}

The following argument independently proves the $q=3$ case using only a finite quartic coefficient test and explicit sums of squares. It does not rely on the full lift.

Now fix a $2+3$ cover. Strip off the product of the five core activities from the sector in which all five core vertices are selected. Its polynomial is
\[
 t^3(c_0+c_1t+c_2t^2).
\]
\begin{theorem}[Sharp constant three]\label{thm:three}
For arbitrary nonnegative exterior activities and every $2+3$ core,
\[
 c_1^2\ge3c_0c_2.
\]
The constant $3$ cannot be increased uniformly.
\end{theorem}

Let $T$ be the weight of exterior-left triples matching all of $B$. Let $C$ be the weight of single exterior-left vertices which, together with the two core rows, match $B$. Let $P_i$ be the weight of exterior-left pairs completable with core row $S_i$, and let $P_U$ use the union row. Set
\[
 O=P_1+P_2-P_U.
\]
The union interpretation gives $0\le O\le\min(P_1,P_2)$. The central left-side inequality is
\begin{equation}\label{eq:leftkey}
 4P_1P_2-2O^2\ge3CT.
\end{equation}

\subsection{A finite exact coefficient certificate}
Encode a nonempty neighborhood in $B$ by its bit mask $1,\ldots,7$. For three nonempty masks, let $M(r,s,t)$ be one exactly when they admit a matching to all three columns. Hall's condition gives the elementary rule: their bitwise union is $7$, and no singleton mask $1,2,4$ appears twice. For general tuples including an empty mask, use ordinary Boolean matchability.

Index left exterior vertices by $i$, with masks $t_i$ and weights $z_i$. Put
\begin{align*}
 c_i&=M(S_1,S_2,t_i),&p_{ij}&=M(S_1,t_i,t_j),&q_{ij}&=M(S_2,t_i,t_j),\\
 o_{ij}&=p_{ij}q_{ij},&\tau_{ijk}&=M(t_i,t_j,t_k).
\end{align*}
Then $C=\sum c_iz_i$, $P_1=\sum_{i<j}p_{ij}z_iz_j$, and similarly for $P_2,O,T$. In the quartic
\[
 I=4P_1P_2-2O^2-3CT
\]
every exponent is at most two. Its coefficient of $z_i^2z_j^2$ is $2o_{ij}\ge0$. Its coefficient of $z_i^2z_jz_k$ for distinct indices is
\begin{equation}\label{eq:repeatcoef}
 4(p_{ij}q_{ik}+p_{ik}q_{ij})-4o_{ij}o_{ik}-3c_i\tau_{ijk}.
\end{equation}
Its coefficient of $z_1z_2z_3z_4$ is
\begin{equation}\label{eq:squarefreecoef}
 4\sum_{ij\mid kl}(p_{ij}q_{kl}+p_{kl}q_{ij})
 -4\sum_{ij\mid kl}o_{ij}o_{kl}
 -3\sum_i c_i\tau_{[4]\setminus\{i\}},
\end{equation}
where each of the three unordered pair partitions is used once.

If the core has matching rank at most one, $C=0$ and \eqref{eq:leftkey} follows from $O^2\le P_1P_2$. Otherwise row exchange and column permutation give exactly the eight representatives in Table~\ref{tab:coeff}. The table records the complete range of possible coefficients in \eqref{eq:repeatcoef} and \eqref{eq:squarefreecoef}. It is obtained by the stated three-row Boolean rule, over $7^3=343$ ordered triples and $\binom{10}{4}=210$ neighborhood multisets per core. The included exact verifier implements an independent permutation matcher and checks all $4424$ substitutions, as well as all $46$ labeled matching-rank-two cores.

\begin{table}[ht]
\centering
\begin{tabular}{ccc}
\toprule
Core masks & Repeated-exponent coefficients & Squarefree coefficients\\
\midrule
$(1,2)$ & $0,1,4$ & $0,2,3$\\
$(1,3)$ & $0,1,4$ & $0,2,3,6,8,12$\\
$(1,6)$ & $0,1,4$ & $-4,0,2,3$\\
$(1,7)$ & $0,1,4$ & $-4,0,2,3,8,12$\\
$(3,3)$ & $0,1,4$ & $0,2,3,6,8,12$\\
$(3,5)$ & $0,1,4$ & $0,2,3,8,12$\\
$(3,7)$ & $0,1,4$ & $0,2,3,8,12$\\
$(7,7)$ & $0,1,4$ & $0,2,3,8,12$\\
\bottomrule
\end{tabular}
\caption{The complete finite coefficient test. In each exceptional row the only four negative neighborhood multisets are listed in \eqref{eq:bad}.}\label{tab:coeff}
\end{table}

Only cores $(1,6)$ and $(1,7)$ have negative coefficients. In both cases the negative neighborhood multisets are exactly
\begin{equation}\label{eq:bad}
 (2,3,4,5),\quad(2,3,5,5),\quad(3,3,4,5),\quad(3,3,5,5),
\end{equation}
and each coefficient is $-4$. Thus every term involving another neighborhood class is nonnegative. It remains to treat the four classes $2,3,4,5$ in these two cores.

\subsection{The four-class sum-of-squares certificate}
For core $(1,6)$ let $a,b,c,d$ now denote the activity sums in classes $2,3,4,5$, and let $B_2,D_2$ be the second elementary moments of classes $3,5$. Put $u=a+b$, $v=c+d$. Exact support counting yields
\begin{align*}
 C&=u+v,& P_1&=uv,\\
 P_2&=(a+c)(b+d)+bd+B_2+D_2,& O&=ad+bc+bd,\\
 T&=v(ab+B_2)+u(cd+D_2).
\end{align*}
Since $0\le B_2\le b^2/2$ and $0\le D_2\le d^2/2$, and $I$ is affine in both moments, it suffices to check the four rectangle corners. Write $I_{00},I_{01},I_{10},I_{11}$ for the lower/lower, lower/upper, upper/lower and upper/upper choices. The following exact identities have only nonnegative summands:
\begin{align*}
 I_{00}={}&2(ad-bc)^2+uv(ab+cd)+(ac+bd)(ad+bc)+2b^2d^2,\\
 I_{01}={}&\tfrac12d^2(a-b)^2+\tfrac12ad(b-d)^2
          +ac(b-d)^2+bc(a-d)^2+R_{01},\\
 R_{01}={}&a^2(bd+cd)
 +a\left(\tfrac12b^2d+bc^2+c^2d+\tfrac12cd^2\right)\\
 &+b^2(2c^2+cd)+b\left(c^2d+\tfrac12cd^2+\tfrac12d^3\right),\\
 I_{11}={}&\tfrac12bd(u-v)^2+\tfrac12(ad+bc+3ac)(b-d)^2\\
 &+(b+d)\left(ac(a+c)+\tfrac12(a^2d+bc^2)\right).
\end{align*}
The expression for $I_{10}$ is obtained from $I_{01}$ by interchanging $(a,b)$ with $(c,d)$. For core $(1,7)$, the restricted $P_2,O$ both increase by $ac$ while $C,T,P_1$ are unchanged. Since $P_1-O=ac$, its quartic is exactly the preceding $I+2a^2c^2$. This proves \eqref{eq:leftkey} in every case. The included symbolic checker verifies all five identities exactly.

\subsection{Passing to the right exterior and proving sharpness}
Return to the right-side notation $a,b,c,d$ of Section 3. Selecting all core vertices gives
\[
 c_0=C,\qquad c_1=aP_2+bP_1+cP_U,\qquad c_2=TQ.
\]
Let $A=a+c$, $B=b+c$, $p=P_1$, $q=P_2$. Then
\[
 c_1=qA+pB-Oc,\qquad Q\le AB-c^2/2.
\]
For $pq>0$, the exact quadratic identity
\begin{align*}
 &(qA+pB-Oc)^2-(4pq-2O^2)(AB-c^2/2)\\
 &\quad=\left(1-\frac{O^2}{2pq}\right)(qA-pB)^2
 +\frac{[O(qA+pB)-2pqc]^2}{2pq}\ge0
\end{align*}
uses only $O^2\le pq$. Together with \eqref{eq:leftkey} it implies $c_1^2\ge3CTQ=3c_0c_2$. If $pq=0$, the same assertion follows by continuity, or directly from $O=0$ and $CT=0$.

For sharpness, take the complete $2\times3$ core, $N$ common left-exterior vertices and $M$ common right-exterior vertices, all with unit exterior activities. Then
\[
 c_0=N,\quad c_1=\binom N2M,\quad c_2=\binom N3\binom M2,
 \qquad
 \frac{c_1^2}{c_0c_2}=3\frac{N-1}{N-2}\frac{M}{M-1}\longrightarrow3.
\]
This completes Theorem~\ref{thm:three}.
